\begin{table}[!t]
\centering
\caption{Principales notations post-A8 du modèle MCAD.}
\scriptsize
\begin{tabularx}{\columnwidth}{@{}p{2.35cm}X@{}}
\toprule
Symbole & Description \\
\midrule
$\mathrm{CKG}$ & Graphe de connaissances contextuel relativement stable \\
$V^{K},A^{K},\tau^{K}$ & Nœuds, arêtes et typage du CKG \\
$O$ & Ensemble fini des contraintes de l'objectif analytique actif \\
$G_O$ & Graphe objectif--évidence \\
$C_O,S_O,E_O$ & Contraintes, supports suffisants et atomes d'évidence de $G_O$ \\
$\operatorname{Supp}(c)$ & Supports suffisants alternatifs de la contrainte $c$ \\
$\operatorname{Req}(s)$ & Atomes d'évidence requis par le support $s$ \\
$\mathcal{E}_t$ & Évidence analytique effectivement acquise par la session \\
$\operatorname{Comp}(c,\mathcal{E}_t)$ & Calculabilité de $c$ à partir de l'évidence acquise \\
$C_{\mathcal{E}}^{\le t}(O)$ & Contraintes calculables cumulativement au temps $t$ \\
$U(\mathcal{E}_t)$ & Frontière résiduelle d'évidence encore utile à l'objectif \\
$P_q$ / CQP & Profil canonique de la requête candidate $q$ \\
$N_E(q)$ & Sur-ensemble sûr de l'évidence potentielle pertinente de $q$ \\
$R_E(q)$ & Évidence potentielle dont la réalisabilité a été établie \\
$A_t$ & Évidence effectivement acquise à l'étape $t$ \\
$\operatorname{ADM}(P_q)$ & Admissibilité sémantique du CQP, hors non-vacuité \\
$\phi_t,\Delta\phi_t$ & Progression et gain immédiat de calculabilité \\
$N_E(q)\cap U(\mathcal{E}_t)$ & Test suffisant utilisé pour l'élagage sûr \\
\bottomrule
\end{tabularx}
\label{tab:notation}
\normalsize
\end{table}
